\documentclass[12pt, letterpaper]{article}

%-------Packages---------
\usepackage{newpxtext,newpxmath} % Palatino-like text & math (modern)
\usepackage{bboldx}
\usepackage{mathtools}
\usepackage{tikz}
\usetikzlibrary{calc,intersections}
\usepackage{tikz-cd}
\usepackage[all,arc]{xy}
\usepackage{graphicx}

\usepackage{enumerate}
\usepackage[dvipsnames]{xcolor}
\usepackage{float}
\usepackage[left=1in,top=1in,right=1in,bottom=1in]{geometry}
\usepackage{fancyhdr}
\usepackage{multicol}
\usepackage{setspace}

\usepackage{dsfont}
\usepackage{mathrsfs}

\usepackage{tcolorbox}
\tcbuselibrary{theorems,skins,breakable}

\usepackage{xparse}
\usepackage{import}
\usepackage[utf8]{inputenc}

\usepackage{hyperref}
\usepackage{cleveref}

\usepackage{xcolor, ifthen, tcolorbox}
\tcbuselibrary{theorems,skins,breakable}
% Theorem System
% The following environments are provided:
%   New Problem:        \begin{problem} ...
%   New Question Header:   \qh (then followed by subproblems)
%   Subproblem:     \begin{subproblem} ...
%   Problem Proof:  \begin{problemproof} ...
%   Proof:          \\begin{pf}{color} ...
%   Remark:         \rmk (sentence), \rmkb (block)

% Define custom colors
\definecolor{thmscol}{HTML}{F4565B} %For Problems
\definecolor{lemscol}{HTML}{FE844C} %For Lemmes
\definecolor{prpscol}{HTML}{00A7EF} %For proposition
\definecolor{corscol}{HTML}{23DAFF} %For corrolaries
\definecolor{rmkscol}{HTML}{6DEEC5} %For remarks
\definecolor{defscol}{HTML}{18C991} %For definitions
\definecolor{exmscol}{HTML}{7DB800} %For examples
\definecolor{clmscol}{HTML}{165c58} %For claims

%Change between dark(black)/light(white) mode
\newcommand{\colormode}{white}
\ifthenelse{\equal{\colormode}{white}}
      {\newcommand{\TextColor}{black}}
      {\newcommand{\TextColor}{white}}
\newcommand{\pbcolormain}{thmscol!80!\colormode}
\newcommand{\pbcolorsecond}{thmscol!38!\colormode}


% ============================
% Problem Counter
% ============================
\newcounter{subproblemcounter}
\newcounter{problemcounter}

\newcommand{\q}{
    \setcounter{subproblemcounter}{0}%
    \stepcounter{problemcounter} % Increment problem counter
    \color{\TextColor}Problem \arabic{problemcounter}: % Display the problem counter
}

\newcommand{\stepsub}{
    \stepcounter{subproblemcounter}%
    \color{\TextColor}(\alph{subproblemcounter})
}
% ============================
% Problem
% ============================
\newtcolorbox{pb}[1]{
  enhanced,
  frame hidden,
  titlerule=0mm,
  toptitle=1mm,
  bottomtitle=1mm,
  fonttitle=\bfseries\large,
  coltitle=black,
  colbacktitle=\pbcolormain,
  colback=\pbcolorsecond,
  title={\q #1},
  coltext=\TextColor
}

\newtcolorbox{spb}[1]{
  enhanced,
  frame hidden,
  titlerule=0mm,
  toptitle=1mm,
  bottomtitle=1mm,
  fonttitle=\bfseries\large,
  coltitle=black,
  colbacktitle=\pbcolormain,
  colback=\pbcolorsecond,
  title={\stepsub #1},
  coltext=\TextColor,
  attach boxed title to top left={xshift=3mm,yshift=-2mm},
  boxed title style={
    colback=\pbcolormain,
    colframe=\pbcolormain,
  }
}

\newtcolorbox{pbh}[1]{
    enhanced,
    frame hidden,
    titlerule=0mm,
    toptitle=1mm,
    bottomtitle=1mm,
    fonttitle=\bfseries\large,
    coltitle=black,
    colbacktitle=\pbcolormain,
    colback=\pbcolorsecond,
    title={\q #1},
    boxrule=0pt,
    interior hidden,
    attach boxed title to top left={xshift=3mm,yshift=-2mm},
    boxed title style={
        colback=\pbcolormain,
        colframe=\pbcolormain,
    }
}

\newenvironment{problem}[1]{%
  \newpage
  \begin{pb}{#1}
    }{
  \end{pb}
}

\newenvironment{subproblem}[1]{%
  \begin{spb}{#1}
    }{
  \end{spb}
}

\newenvironment{problemproof}{%
    \begin{pf}{\pbcolormain}
        }{
    \end{pf}
}

\newcommand{\qh}[1][]{
    \newpage
    \begin{pbh}{#1}\end{pbh} % Display the problem counter
}

% ============================
% Claim
% ============================

\newtcbtheorem[number within=problemcounter]{claim}{Claim}
{
    enhanced,
    frame hidden,
    titlerule=0mm,
    toptitle=1mm,
    bottomtitle=1mm,
    fonttitle=\bfseries\large,
    coltitle=black,
    colbacktitle=clmscol!50!\colormode,
    colback=clmscol!20!\colormode,
}{clm}

\NewDocumentCommand{\clm}{m+m}{
    \begin{claim*}{#1}{}
        #2
    \end{claim*}
}

\newenvironment{clmpf}{
	{\noindent{\it \textbf{Proof.}}
	\setlength{\parindent}{0pt}
	}
	\tcolorbox[blanker,breakable,left=5mm,parbox=false,
    before upper={\parindent15pt},
    after skip=10pt,
	borderline west={1mm}{0pt}{clmscol!50!\colormode}]
}{
    \textcolor{clmscol!50!\colormode}{\hbox{}\nobreak\hfill$\blacksquare$}
    \endtcolorbox
}

\NewDocumentCommand{\clmp}{m+m+m}{
    \begin{claim*}{#1}{}
        #2
    \end{claim*}

    \begin{clmpf}
        #3
    \end{clmpf}
}


% ============================
% Proof
% ============================

\newenvironment{pf}[1]{%
  \def\pfcolor{#1}% Store the color in a macro
  \noindent{\it \textbf{Proof.}}%
  \tcolorbox[blanker,breakable,left=5mm,parbox=false,%
    before upper={\parindent15pt},%
    after skip=10pt,%
    borderline west={1mm}{0pt}{\pfcolor},%
    coltext=\TextColor
  ]%
  \setlength{\parindent}{0pt}%
}{%
  \textcolor{\pfcolor}{\hbox{}\nobreak\hfill$\blacksquare$}%
  \endtcolorbox%
}

\newtcolorbox{solution}{
  blanker,
  breakable,
  left=5mm,
  parbox=false,%
  before upper={\parindent15pt},%
  after skip=10pt,%
  borderline west={1mm}{0pt}{\pbcolormain},%
  coltext=\TextColor
}


% ============================
% Remark
% ============================


\NewDocumentCommand{\rmk}{+m}{
    {\it \color{rmkscol!100!white}#1}
}

\newenvironment{remark}{
    \par
    \vspace{5pt}
    \begin{minipage}{\textwidth}
        {\par\noindent{\textbf{Remark.}}}
        \tcolorbox[blanker,breakable,left=5mm,
        before skip=10pt,after skip=10pt,
        borderline west={1mm}{0pt}{rmkscol!80!\colormode},
        coltext=\TextColor
        ]
}{
        \endtcolorbox
    \end{minipage}
    \vspace{5pt}
}

\NewDocumentCommand{\rmkb}{+m}{
    \begin{remark}
        #1
    \end{remark}
}


% Define a new command for problems
\renewcommand{\headrule}{\color{\TextColor}\hrulefill}
\newcommand{\C}{\mathbb{C}}
\newcommand{\R}{\mathbb{R}}
\newcommand{\Q}{\mathbb{Q}}
\newcommand{\Z}{\mathbb{Z}}
\newcommand{\N}{\mathbb{N}}
\newcommand{\p}{\mathfrak{p}}
\newcommand{\F}{\mathbb{F}}
\newcommand{\contr}{\Rightarrow \Leftarrow}
\newcommand{\s}{\(\,\)}
\renewcommand{\t}[1]{\text{#1}}
\newcommand{\Spec}{\mathrm{Spec}}
\newcommand{\mf}[1]{\mathfrak{#1}}

% Meta data
\setstretch{1.2}

\begin{document}
\color{\TextColor}
\pagecolor{\colormode}
\setlength{\parindent}{0pt}
\pagestyle{fancy}
\fancyhf{}
\fancyhead[LOE]{\color{\TextColor}\slshape{\textbf{Vincent Ferrara}}}
\fancyhead[ROE]{\color{\TextColor}HW4 --- \thepage}
\fancyhead[C]{\color{\TextColor}Math 614}

\begin{problem}{}
    Let $f : R \to S$ be a ring homomorphism. Let $\p \in \Spec(R)$. Let $S_\p$ be the localization of
    $S$ at the multiplicative system $f (R \setminus \p)$. Let $\p S_\p$ be the ideal generated by the image of
    $\p$ under the composition $R \overset{f}{\to}  S \overset{g}{\to} S_\p$, where $g$ is the localization homomorphism. Let
    $h : S \to S_\p \diagup \p S_\p$ be the composition of the localization homomorphism with the quotient
    map.
\end{problem}

\begin{subproblem}{}
    Prove that $\Spec(h)$ induces a homeomorphism
\[\Spec(S_\p \diagup \p S_\p) \to \Spec(f)^{-1}(\p)\]
onto the fiber of the morphism $\Spec(f ) : \Spec(S) \to \Spec(R)$ over $\p$
\end{subproblem}
\begin{problemproof}
    Shown in class $\Spec(g'): \Spec(S_\p) \to D(f(R \setminus \p))$ is a homeomorphism, where $g' = g\mid_{D(f(R \setminus \p))}$ and
    
    $\Spec(\pi): \Spec(S_\p \diagup \p S_\p) \to V(\p S_\p)$ is a homeomorphism where $\pi$ is the quotient map.

    WTS: $\Spec(g')(V(\p S_\p))=\Spec(f)^{-1}(\p)$

    Let $\mf{q} \in \Spec(g')(V(\p S_\p))$ this means there is an 
    $I \in \Spec(S_\p)$ s.t. $\p S_\p \subset I$ and $\Spec(g')(I)=\mf{q}$.

    Since $g(f(p)) \subset \p S_\p \subset I \implies f(p) \subset g^{-1}(I)=q \implies \p \subset f^{-1}(q)$.

    Also, since $g(V(\p S_\p)) \subset D(f(R \setminus \p)) \implies \mf{\q} \cap f(R \setminus \p)= \emptyset$.

    So, let $r \notin \p \implies f(r) \in f(R \setminus \p) \implies f(r) \notin \mf{q} \implies r \notin f^{-1}(\mf{q})$.

    Thus, $f^{-1}(\mf{q}) \subset \p$.

    Thus, $f^{-1}(\mf{q}) \subset \p$, so $\Spec(g')(V(\p S_\p)) \subset \Spec(f)^{-1}(\mf{q})$.

    Let $q \in \Spec(f)^{-1}(\mf{p}) \implies f^{-1}(\mf{q})=\p$

    Let $r \in f(R \setminus \p) \implies \exists a \in R \setminus \p$ s.t. $f(a)=r$. Therefore, $a \notin \p \implies a \notin f^{-1}(\mf{q}) \implies r \notin \mf{q}$.

    Thus, $\mf{q} \cap f(R \setminus \p) = \emptyset$. This implies that $\mf{q}$ extends to $\mf{q} S_\p$ s.t. $g^{-1}(\mf{q} S_\p)=\mf{q}$ due to the correspondence between primes that don't meet $f(R \setminus \p)$ and primes in $S_\p$.

    Also, since $\p = f^{-1}(\mf{q}) \implies f(\p) \subset \mf{q} \implies g(f(\p))=\mf{q} S_\p \implies \p S_\p \subset \mf{q} S_\p$ since $\p S_\p$ is generated by $g(f(\p))$.

    Thus, $g(\mf{q}) \in V(\p S_\p) \implies \mf{q} \in \Spec(g')(V(\p S_\p))$.

    Therefore, $\Spec(g')(V(\p S_\p))=\Spec(f)^{-1}(\p)$.

    This implies that $\Spec(h)=\Spec(g') \circ \Spec(\pi) : \Spec(S_\p \diagup \p S_\p) \to \Spec(g')(V(\p S_\p))=\Spec(f)^{-1}(\p)$ is a homomorphism since it is a composition of homeomorphisms.
\end{problemproof}
\stepcounter{subproblemcounter}
\begin{subproblem}{}
    Explain why $S_\p \diagup \p S_\p$ is naturally a $\kappa(\p)$-algebra, where $\kappa(\p)$ is the residue field of $R$
at $\p$.
\end{subproblem}
\begin{problemproof}
    $f : R \to S$ extends to $f_\p : R_\p \to S_\p$, and let $\pi : S_\p \to S_\p \diagup \p S_\p$ be the quotient map.

    Define $\varphi : R_\p \to S_\p \diagup \p S_\p$ s.t. $\varphi=\pi \circ f_\p$.

    Let $a/s \in \p R_\p \implies a \in \p$ and $s \in R \setminus \p$.

    Thus, $\varphi(a / s)=\pi(f(a)/f(s))=\pi(f(a)/1)\pi(1/f(s))$. Since $f(a)/1 \in \p S_\p$ this implies that $\pi(f(a)/1)=0$.

    Thus, $\varphi(a/s)$, so $p R_\p \subset \ker(\varphi)$. Therefore, by the universal property of quotients there exists a unique
    $\bar{\varphi} : R_\p \diagup \p R_\p \to S_\p \diagup \p S_\p$ ring homomorphism s.t. $\varphi=\bar{\varphi} \circ \pi_R$, where $\pi_R$ is the quotient map for $R$.

    Thus, $S_\p \diagup \p S_\p$ is a $\kappa(\p)$-algebra.
\end{problemproof}

\begin{subproblem}{}
    Prove that if $f : R \to S$ is finite, then $S_\p \diagup \p S_\p$ is a finite $\kappa(\p)$-algebra.
\end{subproblem}
\begin{problemproof}
    Assume $f: R \to S$ is finite. This implies $S=f(R)x_1 \oplus \dots \oplus f(R)x_n$ for $x_i \in S$.

    Consider $\bar{\varphi} : \kappa(\p) \to S_\p \diagup \p S_\p$ the construction from the last part.

    Let $\overline{a/s} \in S_p \diagup \p S_\p$, and let $a/s \in S_\p$ s.t. $\pi(a/s)=\overline{a/s}$, and since $s \in f(R \setminus \p) \implies s=f(b)$ for $b \in R \setminus \p$.

    Since $a \in S$ this implies that $a=\sum_{i=1}^nf(r_i)x_i$ for $r_i \in R$.

    Therefore, $\pi(a/s)=\pi(a/1)\pi(1/s)=\sum_{i=1}^n \pi(x_i)\pi(f(r_i)/f(b))$.

    Since $f(r_i)/f(b)=f_\p(r_i/b) \implies \pi(f(r_i)/f(b))=\pi(f_\p(r_i/b))=\varphi(r_i/b)$.

    Thus, $\pi(a/s)=\sum_{i=1}^n\pi(x_i/1)\varphi(r_i/b)=\sum_{i=1}^n\pi(x_i/1)\bar{\varphi}(\pi_R(r_i/b))$.

    Thus, $S_\p \diagup \p S_\p = \bigoplus_{i=1}^n \bar{\varphi}(R_\p \diagup \p R_\p) \pi(x_i/1)$.

    Thus, $S_\p \diagup \p S_\p$ is a finite $\kappa(\p)$-algebra.
\end{problemproof}

\begin{subproblem}{}
    Prove that if $A$ is a finite $k$-algebra where $k$ is a field, then $\Spec(A)$ is a finite set.
\end{subproblem}
\clmp{}{
    If $R$ is a noetherian ring then $R^n$ is a noetherian $R$-module.
}{
    \setlength{\parindent}{0pt}
    Base Case $k=1$: $R$ is noetherian.

    Inductive Case: Assume $R^k$ is noetherian for $k \leq n$.

    Consider $R^{n+1}$.

    Then since $0 \to R \to R^{n+1} \to R^n \to 0$ is a short exact sequence. This implies that $R^{n+1}$ is noetherian since $R$ and $R^{n}$ is noetherian by inductive hypothesis.

    Thus, $R^n$ is noetherian for all $n \in \N$.
}
\clmp{}{
    If $R$ is a noetherian ring and $M$ is a finitely generated $R$-module then $M$ is noetherian.
}{
    \setlength{\parindent}{0pt}    
    Assume $R$ is a noetherian ring and $M$ is a finitely generated $R$-module.

    Since $M$ is finitely generated this implies there exists $x_1,\dots,x_n \in M$ s.t. $M=\bigoplus_{i=1}^n Rx_i$.

    Thus, there is a surjective map $\phi : R^n \to M$ s.t. $\varphi(r_1,\dots,r_n)=\sum_{i=1}^n r_ix_i$.

    This implies that $R^n \diagup \ker(\phi) \cong M$ shown above since $R^n$ is noetherian and $M$ is isomorphic to a quotient of $R^n$ this implies $M$ is noetherian.
}
\clmp{}{
    If $A$ is a ring and is noetherian as an $R$-module then $A$ is noetherian.
}{
    \setlength{\parindent}{0pt}
    Assume $A$ is a ring and is noetherian as an $R$-module by the map $\varphi : R \to A$.

    Let $I$ be an ideal in $A$. Since $I$ is a submodule of $A$ this implies that $I$ is
    finitely generated as an $R$-module. Thus, $I=\bigoplus_{i=1}^n \varphi(R)x_i$ for $x_i \in I$.

    Consider $J=(x_1,\dots,x_n)$.

    First $J \subset I$ since $I$ is an ideal and contains $x_1,\dots,x_n$.

    Let $y \in I$ then $y=\sum_{i=1}^n \varphi(r_i)x_i$ for $r_i \in R$.

    This implies that $y \in J$ since $\varphi(r_i) \in A$.

    Thus, $I=J$ and implies $I$ is a finitely generated. Therefore, $A$ is a noetherian ring since every ideal is finitely generated.
}
\begin{problemproof}
    Assume $A$ is a finite $k$-algebra where $k$ is a field.

    This implies $A$ is a finite dim $k$-vector space with dimension $n$.

    This implies there is a map $f : k \to A$ s.t. $f$ is finite.

    Let $\p \subset A$ be a prime ideal.

    Let $x \notin \p$. This implies $x \neq 0$ thus $\{1,\dots,x^n\}$ is linear dependent.

    Thus, there exists $c_i \in f(k)$ s.t. $c_0 + c_1x+\dots+c_nx^n=0$ s.t. some $c_i \neq 0$.

    Now map this relation to the quotient $A \diagup \p$. Since $x \notin \p \implies \bar{x} \neq 0$.

    Thus, $\bar{c_0} + \bar{c_1}\bar{x} + \dots + \bar{c_n}\bar{x}^n = 0$.

    If $c_i \neq 0$ then $c_i$ is a unit since it is in $f(k)$ and ring homomorphism preserve units. Thus, if $c_i \neq 0 \implies c_i \notin \p$ thus $\bar{c_i} \neq 0$.

    Now if $\bar{c_0} = 0$ then $\sum_{i=1}^n \bar{c_i}\bar{x}^n=0 \implies \bar{x} \sum_{i=1}^n \bar{c_i}\bar{x}^{i-1}=0$, but $A\diagup \p$ is an integral domain, so this is a contradiction.

    Therefore, $\bar{c_0} \neq 0$.

    Thus, $-\bar{c_0}(\sum_{i=1}^n \bar{c_1}\bar{x}^{i-1}) \bar{x} = 1 \implies \bar{x}$ is a unit. Thus, $A \diagup \p$ is a field, so $\p$ is maximal.

    Thus, all prime ideals in $A$ are maximal.
    
    Let $\p, \mf{q} \in \Spec(A)$ s.t. $\p \subset \mf{q}$ since $\p$ is maximal this implies that $\p = \mf{q}$ thus all primes are minimal.

    Also, since $k$ is noetherian this implies that $A$ is a noetherian ring. Thus, shown in class this implies that $A$ is a finite number of minimal primes, but since all primes are minimal $\Spec(A)$ is finite.
\end{problemproof}

\begin{subproblem}{}
    Prove that if $f : R \to S$ is finite, then the fibers of $\Spec(f) : \Spec(S) \to \Spec(R)$ are
finite sets.
\end{subproblem}
\begin{problemproof}
    Assume $f : R \to S$ is finite then for $\p \in \Spec(R)$ we know that $S_\p \diagup \p S_\p$ is a finite $\kappa(\p)$-algebra by (c).
    
    Then by (d) since $\kappa(\p)$ is a field we know that $S_\p \diagup \p S_\p$ has finite primes.

    Therefore, by (a) since $\Spec(f)^{-1}(\p)$ is homeomorphic to $\Spec(S_\p \diagup \p S_\p)$ this implies the fibers of $\Spec(f)$ are finite.
\end{problemproof}
\setcounter{problemcounter}{1}
\begin{problem}{}
    
\end{problem}
\begin{subproblem}{}
    Check that a ring homomorphism $f : R \to S$ satisfies going up if and only if 
    specializations lift along the map $\Spec(f ) : \Spec(S) \to \Spec(R)$.
\end{subproblem}
\begin{problemproof}
    $(\implies)$

    Assume $f: R \to S$ satisfies going up.

    Let $\p, \p' \in \Spec(R)$, and $\mf{q} \in \Spec(S)$ s.t. $\p' \in \overline{\{\p\}}=V(\p)$ and $f^{-1}(\mf{q})=\p$.

    Since $\p' \in V(\p) \implies \p \subset \p'$ thus since $f$ satisfies going up this implies that there is a $\mf{q}' \in \Spec(S)$ s.t. $f^{-1}(\mf{q}')=\p'$ and $\mf{q} \subset \mf{q}'$.

    Thus, $\mf{q}' \in V(\mf{q})$, so this implies that specializations lift along $\Spec(f)$.

    $(\impliedby)$

    Assume specializations lift along the map $\Spec(f)$.

    Let $\p, \p' \in \Spec(R)$, and $\mf{q} \in \Spec(S)$ s.t. $\p \subset \p'$ and $f^{-1}(\mf{q})=\p$. Since $\p \subset \p' \implies \p' \in V(\p)$. Thus, by assumption
    
    there exists a $\mf{q}' \in \Spec(S)$ s.t. $\mf{q}' \in V(\mf{q})$ and $f^{-1}(\mf{q}')=\p'$. Thus, $\mf{q} \subset \mf{q}'$, so $f$ satisfies going up.
\end{problemproof}

\begin{subproblem}{}
    Let $f : R \to S$ be a ring homomorphism and let $T \subset Spec(R)$ be the image of the
map $Spec(f ) : Spec(S) \to Spec(R)$. Assume that $T$ satisfies the following property: if
$\p \in Spec(R)$ is any point which is a specialization of a point of $T$, then $\p \in T$. Prove
that $T$ is closed.
\end{subproblem}
\begin{problemproof}
    Let $\p \in \overline{T}$.

    Assume $S_\p = 0$. This implies that there is an $r \in R \setminus \p$ s.t. $f(r)=0$. 
    
    This implies that $f(r) \in \sqrt{(0)} \implies f(r) \in \mf{q}$ for $\mf{q} \in \Spec(S)$. Thus, $r \in f^{-1}(\mf{q})$ for all $\mf{q} \in \Spec(S)$. Thus, $T \subset V(r) \implies \overline{T} \subset V(r)$.

    Thus, $r \in \p$, but this is a contradiction since $r \in R \setminus \p$.

    Thus, $S_\p \neq 0$.

    Let $\mf{q} S_\p \in \Spec(S_p)$ this corresponds to the prime $\mf{q} \in \Spec(S)$ s.t. $\mf{q} \cap f(R \setminus \p) = \emptyset$.

    Let $\p' = f^{-1}(q)$.

    Let $x \in \p' \implies f(x) \in \mf{q} \implies x \notin f(R \setminus \p) \implies x \in \p \implies \p' \subset \p$.
    
    Thus, $\p \in V(\p')$, so $\p'$ is a specialization of $\p$. Since $\p' \in T \implies \p \in T$. Thus, $\overline{T} = T$.
\end{problemproof}

\begin{subproblem}{}
    For a ring homomorphism $f : R \to S$, prove that $f$ satisfies going up if and only if the
map $\Spec(f ) : \Spec(S) \to \Spec(R)$ is closed.
\end{subproblem}
\begin{problemproof}
    $( \implies )$

    Assume $f$ satisfies going up.

    Let $I$ be an ideal of $S$. Consider $f^{-1}(V(I))=T$.

    Let $\p' \in \Spec(R)$ which is a specialization of $\p \in T$. Thus, $\p' \in V(\p) \implies \p \subset \p'$, since $\p \in T$ this implies there exists a $\mf{q} \in V(I)$ s.t. $f^{-1}(\mf{q})=\p$.

    Since $f$ staisfies going up this implies that there is a $\mf{q}' \in \Spec(S)$ s.t. $I \subset \mf{q} \subset \mf{q}'$ and $f^{-1}(\mf{q}')=\p'$.

    Thus, $\mf{q}' \in V(I) \implies \p' \in T$. Thus, $T$ is closed by (b), so $\Spec(f)$ is a closed map.

    $(\impliedby)$

    Assume $\Spec(f) : \Spec(S) \to \Spec(R)$ is closed.

    Let $\p,\p' \in \Spec(R)$ and $\mf{q} \in \Spec(S)$ s.t. $\p \subset \p'$ and $f^{-1}(\mf{q})=\p$.

    Consider $f^{-1}(V(\mf{q}))=V(I)$ for an ideal $I \subset R$.

    Since $I \subset \p \subset \p' \implies \p' \in V(I)$ thus there exists a $\mf{q}' \in V(\mf{q})$ s.t. $f^{-1}(\mf{q}')=\p'$, and since $\mf{q}' \in V(\mf{q}) \implies \mf{q} \subset \mf{q}'$. Thus, $f$ satisfies going up.
\end{problemproof}

\begin{problem}{}

\end{problem}
\begin{subproblem}{}
    Let $\mathrm{Alg}_k^{\mathrm{red,ft}}$ be the full subcategory of $\mathrm{Alg}_k$ consisting of reduced
    finite type $k$-algebras. Explain why $\mathcal{O}$ can also be regarded as a contravariant functor
    \[\mathcal{O} : \mathrm{Aff}_k \to \mathrm{Alg}_k^{\mathrm{red,ft}}\]
\end{subproblem} 
\begin{problemproof}
    We can view the restriction of the codomain as a functor since if $X$ is an affine variety, then $\mathcal{O}(X) \cong k[x_1,\dots,x_n] \diagup I(X)$.

    This implies that $\mathcal{O}(X)$ is finite type since it is finitely generated as a $k$-algebra.

    Also, since $I(X)$ is radical this implies that the $(0)$ ideal in $\mathcal{O}(X)$ is radical thus, $\sqrt{(0)}=(0)$. Shown in class this implies that $\mathcal{O}(X)$ is reduced.

    All arrows still land inside the subcategory since for a regular map $f : X \to Y$ we get the map $f^*: \mathcal{O}(Y) \to \mathcal{O}(X)$ which is between reduced finite type $k$-algebras.
    Therefore, it can be regarded as a contravariant functor.
\end{problemproof}
\begin{subproblem}{}
    Let $(\mathrm{Aff}_k)^\mathrm{op}$ denote the opposite category of $\mathrm{Aff}_k$. Prove that the functor $\mathcal{O}$ gives an
equivalence of categories $(\mathrm{Aff}_k)^\mathrm{op} \simeq \mathrm{Alg}^\mathrm{red,ft}_k$.
\end{subproblem}
\begin{problemproof}
    Shown in HW1 we know that $\mathcal{O}$ is fully faithful on $\mathrm{Alg}_k$ since
    
    $\mathrm{Hom}_{(\mathrm{Aff}_k)^\mathrm{op}}(Y,X) = \mathrm{Hom}_{\mathrm{Aff}_k}(X,Y)$

    Also since shown in part (a) that $\mathcal{O}$ can be viewed as a contravariant functor on the reduced finite $k$-algebras this implies that $\mathcal{O}$ is fully faithful on reduced finite type $k$-algebras.

    Let $A \in \mathrm{Alg}_k^{\mathrm{red,ft}}$. This implies that $A \cong k[x_1,\dots,x_n] \diagup I$ for some ideal $I$. Since $A$ is reduced this implies that $I$ is radical.

    Thus, define $X := V(I)$. 
    
    Since $I$ is radical, the coordinate ring of $X$ is $\mathcal{O} = k[x_1,\dots,x_n] \diagup I \cong A$.

    Therefore, $\mathcal{O}$ is essentially surjective.

    Thus, $\mathcal{O}$ gives an equivalence of categories.
\end{problemproof}
    
\begin{problem}{}
    Let $R$ be a finite type $k$-algebra, where $k$ is a field.
\end{problem}
\begin{subproblem}{}
Prove that any radical ideal $I \subset R$ is the intersection of all maximal ideals in $R$ which
contain $I$.    
\end{subproblem}
\begin{problemproof}
    Assume $I$ is a radical ideal of $R$. Let $A = R \diagup I$, which is still a finite type $k$-algebra.

    This is because since $R$ is finite type this implies there is a surjection from $k[x_1,\dots,x_n] \twoheadrightarrow R$ then we can compose it with the quotient map which is also surjective, so we have a surjection from $k[x_1,\dots,x_n] \to A$.

    Thus, $A$ is also finite type $k$-algebra.

    Let $m$ be in the intersection of all maximal ideals of $A$.

    Assume for contradiction that $m$ there is a prime $\p \subset A$ s.t. $m \notin \p$

    Let $B=A \diagup \p$, so $B$ is a domain, and also finite-type since a quotient.

    Also, this implies that $m \neq 0$ in $B$. Just replace $m$ with $\bar{m}$ when looking at $B$.

    Then $B_m$ is also a domain and also finite type over $k$ since $B_m \cong B[x] \diagup (mx-1)$, which implies there is a surjection from $k[x_1,\dots,x_n,x] \to B_m$.

    Since $B_m$ is a domain this implies $(0)$ is prime thus we have a maximal ideal $\mf{n}$, so $B_m \diagup \mf{n}$ is a field.

    Shown in class this implies $B_m \diagup \mf{n}$ is a finite over $k$.

    Let $J$ be the kernel of the quotient map from $B_m$ to $B_m \diagup \mf{n}$ composed with the localization map from $B$ to $B_m$.

    This implies that $B \diagup J \to B_m \diagup \mf{n}$ is an injection thus, $B \diagup J$ is a $k$-subspace of a finite dim $k$ vector space.

    Thus, $B \diagup J$ is also finite over $k$.
    
    Also, this implies that $B \diagup J$ is a domain since it is a subring of a domain, so $J$ is prime.

    Consider $f : B \diagup J \to B \diagup J$ s.t. $f(x)=ax$ for $a \neq 0$.

    Let $x,y \in B \diagup J$ then $f(x+y)=a(x+y)=ax+ay=f(x)+f(y)$.

    Let $c \in k$ then $f(cx)=acx=cax=cf(x)$. Thus, $f$ is $k$-linear.

    Since $B \diagup J$ is a domain this implies that $f$ is injective.

    By rank nullity theorem this implies that $f$ is surjective and thus an isomorphism. This implies there is a $b \in B \diagup J$ s.t. $ab=1$.

    This implies every nonzero element is invertible thus $B \diagup J$ is a field.

    Thus, $J$ is a maximal ideal in $B$. Since $J$ is the kernel of the map defined above this implies that it is the contraction of $\mf{n}$ the maximal ideal in $B_m$ which we localized at.
    
    Thus, $J$ is disjoint from $\{1,f,f^2,\dots\}$.

    This implies $m \notin J$.

    Now looking at $A$ we know that the maximal ideals of $B$ correspond to maximal ideals of $A$ that contain $\p$.

    This is because if $M$ is a maximal ideal in $B$ then $B \diagup M$ is a field but since $B \diagup M = (A \diagup I) \diagup (\pi^{-1}(M) \diagup I)$ by the third isomorphism theorem we have that this is isomorphic to $A \diagup \pi^{-1}(M)$ which implies $\pi^{-1}(M)$ is also maximal.

    This implies that $m$ is not in the lifted ideal $J$ in $A$ thus this is a contradiction. Therefore, $m$ is in all prime ideals in $A$.

    Now lifting again this implies that $\bigcap_{I \subset \mf{n} \in \mathrm{mSpec}(R)} \mf{n} = \sqrt{I}=I$.
\end{problemproof}
\begin{subproblem}{}
    Prove that for any closed subset $Z \subset \Spec(R)$, the intersection $\Z \cap \mathrm{mSpec}(R)$ is dense
in $Z$, where $\mathrm{mSpec}(R)$ is the maximal spectrum of $R$.
\end{subproblem}
\begin{problemproof}
    Let $Z=V(I)=V(\sqrt{I})$ for some ideal $I$ in $R$. Assume that $V(I)$ is not empty since if it is then it is trivially dense in itself.

    Consider $B=V(\sqrt{I}) \cap \mathrm{mSpec}(R)$.

    Let $A=D(f) \cap V(\sqrt{I})$ be a basic nonemtpy open set in subpsace topology on $Z$ s.t. $f \in R$.

    Assume for contradiction that $A \cap B=\{\p \in \mathrm{mSpec}(R) \mid \sqrt{I} \subset \p \t{ and } f \notin \p\}$ is empty.
    
    This implies that $f$ is in all maximal ideals that contain $\sqrt{I}$.

    This implies that $A$ is empty since $A$ since if $\p \in A$ then $\sqrt{I} \subset \p$ and $f \notin \p$. Now let $\mf{n}$ be a maximal ideal that contains $\p$.

    This is a contradiction since $f \in \mf{m}$. Thus, $A$ being empty is a contradiction, therefore $A \cap B \neq \emptyset$.

    This implies that $\overline{B} = Z$.
\end{problemproof}

\begin{problem}{}
    For each of the following rings $R$, find the irreducible components of $\Spec(R)$.
\end{problem}
\begin{subproblem}{}
    $R=\C[x,y] \diagup (x^3+x-x^2y-y)$.
\end{subproblem}
\begin{problemproof}
    $(x^3+x-x^2y-y)=(x^2(x-y)+(x-y))=(x-y)(x^2+1)=(x-y)(x-i)(x+i)$

    Thus, $\Spec(R)=V(\bar{x}-\bar{y}) \cup V(\bar{x}-\bar{i}) \cup V(\bar{x}+\bar{i})$.

    These are the irreducible since $V(\bar{x}-\bar{y}) \cong \Spec(\C[t])$, $V(\bar{x}-\bar{i}) \cong \Spec(\C[y])$, and $V(\bar{x}+\bar{i}) \cong \Spec(\C[y])$,
    which are all irreducible since $\C[t]$ is a domain so $\Spec(\C[t])$ is irreducible.

    Therefore, since each are distinct and not contained in the union of the other two, shown in class this implies that these are the irreducible components.
\end{problemproof}
\begin{subproblem}{}
    $R=\Z[x,y,z] \diagup (xy,xz,yz)$
\end{subproblem}
\begin{problemproof}
    Define $\varphi : \Z[x,y,z] \to \Z[x] \times \Z[y] \times \Z[z]$ s.t. $f(x,y,z) \to (f(x,0,0), f(0,y,0), f(0,0,z))$.

    Surjective: Given $(g(x),h(y),k(z))$ we have $\varphi(g(x)+h(y)+k(z))=(g,h,k)$.

    Kernel: If $f \in \ker(\varphi)$ this implies that $f$ has no pure x,y,z only terms. Thus $f$ has a multiple of one of $xy,xz,yz$. Thus, $\ker(\varphi) = (xy,xz,yz)$.

    Let $f \in \ker(\varphi)$ then $\pi_1 \circ \varphi(f)=0, \pi_2 \circ \varphi(f)=0,$ and $\pi_3 \circ \varphi(f)=0$. Thus, $f \in \bigcap_{i=1}^3 \ker(\pi_i \circ \varphi)$.

    Let $f \in \bigcap_{i=1}^3 \ker(\pi_j \circ \varphi)$. Thus, $\varphi(f)=(0,0,0)$ so $f \in \ker(\varphi)$.

    Thus, $\ker(\varphi) = \bigcap_{i=1}^3 \ker(\pi_j \circ \varphi)$. This implies that $I=\ker(\varphi)= (y,z) \cap (x,z) \cap (x,y)$.

    Each of $(y,z),(x,z),(x,y)$ is prime since $\Z[x,y,z] \diagup \ker(\pi_i \circ \varphi) \cong \Z[t]$ and $\Z[t]$ is an integral domain.

    Since $\sqrt{I} = \sqrt{(y,z) \cap (x,z) \cap (x,y)} = \sqrt{(y,z)} \cap \sqrt{(x,z)} \cap \sqrt{(y,z)} = (y,z) \cap (x,z) \cap (y,z)=I$. This implies $I$ is radical.

    Thus, $V(I) = V((y,z) \cap (x,z) \cap (x,y))=V(y,z) \cup V(x,z) \cup V(x,y)$, also since we know that $V(\ker(\pi_i \circ \varphi)) \cong \Spec(\Z[t])$ this implies that these are irreducible since $\Z[t]$ is a domain.

    Consider $V(y,z) \cup V(x,y) = V((y,z) \cap (x,y))=V((y,z)(x,y))=V(xy,y^2,xz,yz)$. Also, by same proof above we know that $(xy,y^2,xz,yz)$ is radical. Now consider $V(x,z)$, since there is no way to get a pure $x,z$ term from $(xy,y^2,xz,yz)$ this implies that $V(x,z) \nsubseteq V(x,y) \cup V(y,z)$.

    Do a similar proof to show for all union of two does not contain the third.

    Thus, shown in class since these are irreducible and closed and no union of two sets contains the third set. This implies that these are irreducible components.

    Since $V(I) \cong R$ this implies that $\Spec(R) = V(\overline{y,z}) \cup V(\overline{x,z}) \cup V(\overline{x,y})$. Therefore, these are the irreducible components of $R$.
\end{problemproof}

\begin{problem}{}
    Let $M$ be a noetherian module over a ring $R$. Let $f : M \to M$ be a surjective $R$-module
homomorphism. Prove that $f$ is an isomorphism.
\end{problem}
\begin{problemproof}
    Assume $M$ is a noetherian $R$-module, and $f : M \to M$ is a surjective homomorphism.

    Consider $\ker(f) \subset \ker(f^2) \subset \dots$ since $M$ is noetherian this implies there exists a $N \in \N$ s.t. for all $n \geq N$ that $\ker(f^N)=\ker(f^n)$.

    Let $x \in \ker(f)$. This implies since $f$ is surjective that $y \in f^N(M)$ s.t. $f^N(y)=x \implies f^{N+1}(y)=0$, but since $\ker(f^N)=\ker(f^{N+1})$ this implies that $y \in \ker(f^N)$, thus $x=0$.

    Therefore, $f$ is injective, and thus a isomorphism.
\end{problemproof}

\begin{problem}{}
    For each of the following, either provide an example or prove that none exists.
\end{problem}
\begin{subproblem}{}
    A subring $R$ of a field $K$ such that $R$ is not noetherian
\end{subproblem}
\begin{problemproof}
    Consider $K=\C(x_1,x_2,\dots)$ and $R=\C[x_1,x_2,\dots]$. $K$ is a field and $R$ is a subring that is not noetherian. $R$ being not noetherian was shown in class.
\end{problemproof}

\begin{subproblem}{}
    A field $k$, a finite type $k$-algebra $R$, and a sub-$k$-algebra $S \subset R$ such that $S$ is not
noetherian.
\end{subproblem}
\begin{problemproof}
    Consider $R=k[x,y]$ this is a finite type $k$-algebra. 
    
    Now consider $S=k[x,xy,xy^2,xy^3,\dots]$ which is a sub-$k$-algebra of $R$.

    Consider the ideal $I=(x,xy,xy^2,\dots)$

    Assume for contradiction this is finitely generated then $I=(xy^{i_1},\dots,xy^{i_n})$ and order the genenrators s.t. $xy^{i_j} < xy^{i_{j+1}}$.

    Consider $xy^{i_n + 1}$ since $I$ does not have any pure $y$ terms this implies there is no way to get $xy^{i_n + 1}$ multiplying also by an $x$.

    Thus, $I$ is not finitely generated, so $S$ is not noetherian.
\end{problemproof}
\end{document}